\documentclass[conference]{IEEEtran}

\usepackage{amsmath,amssymb}
\usepackage{booktabs}
\usepackage{graphicx}
\usepackage{url}

\begin{document}

\title{Reinforcement Learning for Drone Navigation in a 3D Environment}

\author{
\IEEEauthorblockN{Artem Davydenko}
\IEEEauthorblockA{
Technical University of Kosice\\
Kosice, Slovakia\\
artem.davydenko@student.tuke.sk}
\and
\IEEEauthorblockN{Vladyslav Kalashnyk}
\IEEEauthorblockA{
Technical University of Kosice\\
Kosice, Slovakia\\
Vladyslav.kalashnyk@student.tuke.sk}
}

\maketitle

\begin{abstract}
This technical report compares three reinforcement learning algorithms for
controlling a point-like drone in a bounded three-dimensional environment with
obstacles. The implemented agents use Deep Q-Networks (DQN), Proximal Policy
Optimization (PPO), and Soft Actor-Critic (SAC). The task is to navigate from a
fixed start position to a fixed goal while avoiding collisions. The comparison
focuses on the average time to reach the target, training collisions, success
rate, and trajectory quality relative to an A* baseline. The results show that
PPO and DQN solve the task on two of three tested seeds, while SAC does not
converge under the selected training budget.
\end{abstract}

\begin{IEEEkeywords}
reinforcement learning, drone navigation, DQN, PPO, SAC, A*, gymnasium
\end{IEEEkeywords}

\section{Introduction}

Autonomous drone navigation is a useful reinforcement learning (RL) benchmark
because it combines sequential decision making, spatial reasoning, and safety
constraints. Related drone-racing simulators, such as AirSim Drone Racing Lab,
provide photo-realistic environments for autonomy and learning research
\cite{madaan2020airsim}. In this project, the problem is simplified to a
controlled 3D navigation task: a point-like drone must move from a predefined
start position to a predefined target in a bounded cube while avoiding static
box obstacles.

The goal of the assignment is to implement and compare three RL algorithms:
DQN \cite{mnih2015dqn}, PPO \cite{schulman2017ppo}, and SAC
\cite{haarnoja2018sac}. The comparison is based on training-time collisions,
evaluation success rate, average number of steps required to reach the target,
and trajectory quality compared with an optimal reference path. The
implementation is based on \texttt{gymnasium} \cite{gymnasium} and
\texttt{stable-baselines3} \cite{stablebaselines3}.

\section{Methodology}

\subsection{Environment}

The environment is implemented as a \texttt{gymnasium.Env} \cite{gymnasium}.
The world is a $10 \times 10 \times 10$ cube with coordinates in $[-5, 5]^3$.
The start position is $[-4, -4, -4]$, and the goal is $[4, 4, 4]$. An episode
is successful when the Euclidean distance to the goal is at most $0.75$. Each
episode is limited to 150 steps.

Two deterministic axis-aligned box obstacles are placed in the environment.
A collision occurs when the proposed drone position is outside the world
boundary or inside an obstacle. Collisions terminate the episode immediately.

\subsection{State and Action Spaces}

The observation is a 9-dimensional continuous vector:
\begin{equation}
s = [p / 10,\; g / 10,\; (o_{\mathrm{near}} - p) / 10],
\end{equation}
where $p$ is the current drone position, $g$ is the goal position, and
$o_{\mathrm{near}}$ is the center of the nearest obstacle.

DQN uses a discrete action space with six unit moves:
\[
\{+x, -x, +y, -y, +z, -z\}.
\]
PPO and SAC use a continuous action space \texttt{Box(3)} with values clipped
to $[-1, 1]^3$. If the vector norm is greater than one, the action is
normalized before applying the step.

\subsection{Reward Function}

The reward encourages progress toward the target and penalizes unsafe or slow
behavior:
\begin{equation}
r_t = 5 \Delta d - 0.1 + r_{\mathrm{terminal}},
\end{equation}
where $\Delta d$ is the reduction in distance to the goal between consecutive
steps. The terminal reward is:
\[
r_{\mathrm{terminal}} =
\begin{cases}
+100, & \text{goal reached},\\
-100, & \text{collision},\\
-10, & \text{timeout},\\
0, & \text{otherwise}.
\end{cases}
\]

\subsection{Episode Termination and Metrics}

Each episode can finish in one of three mutually exclusive outcomes. A
successful episode ends when the drone reaches the goal threshold. A collision
episode ends when the next proposed position intersects an obstacle or leaves
the world bounds. A timeout episode ends when the maximum number of steps is
reached without success or collision.

During training, a custom callback records the number of completed episodes,
collisions, successful goal reaches, and timeouts. This is important because
the assignment requires the number of collisions during training, not only
collisions during final evaluation. During evaluation, the same terminal events
are recorded together with total reward, number of steps, final distance to the
goal, path length, and the full trajectory of each episode.

\subsection{Optimal Path Baseline}

The baseline path is computed using A* \cite{hart1968astar} on a 6-neighbour
3D grid with step size 1. This action model matches the DQN action space. For
the fixed scene, the straight-line start-to-goal distance is approximately
$13.86$, while the A* path length is 24.0 with 24 steps. PPO and SAC can
produce shorter Euclidean paths than this grid baseline because their actions
are continuous.

\section{Experimental Setup}

Each algorithm was trained on three random seeds: 0, 1, and 2. DQN and PPO
were trained for 200,000 environment steps. SAC was trained for 30,000 steps
because it was substantially slower on CPU in this environment. All algorithms
used an MLP policy from \texttt{stable-baselines3} \cite{stablebaselines3}.
The hyperparameters were chosen separately per algorithm and kept fixed across
seeds.

Evaluation was performed over 30 episodes per seed. Two evaluation modes were
used: deterministic policy evaluation and stochastic policy evaluation. The
reported metrics are averaged over the three seeds.

\subsection{Algorithm Configuration}

All three agents use the same 9-dimensional observation vector but differ in
their action representation and learning objective. DQN is used only with the
discrete six-action space because value-based DQN is designed for discrete
actions. PPO and SAC are used with the continuous action space, which allows
the learned policy to choose diagonal motion vectors.

DQN uses an experience replay buffer and an epsilon-greedy exploration schedule.
PPO is trained on batches of on-policy rollouts and optimizes a clipped policy
objective. SAC is an off-policy maximum-entropy actor-critic algorithm. In this
project, no extensive hyperparameter search was performed; instead, reasonable
per-algorithm settings were selected and kept fixed for all seeds. This keeps
the comparison reproducible and prevents tuning one algorithm more aggressively
than the others.

Table~\ref{tab:hyperparams} summarizes the most important hyperparameters used
in the experiment. The values were selected to keep the experiment small enough
for CPU execution while still allowing DQN and PPO to learn meaningful
navigation behavior.

\begin{table}[htbp]
\caption{Main training hyperparameters.}
\label{tab:hyperparams}
\centering
\begin{tabular}{lrrr}
\toprule
Parameter & DQN & PPO & SAC \\
\midrule
Learning rate & $10^{-3}$ & $3\cdot10^{-4}$ & $3\cdot10^{-4}$ \\
Batch size & 128 & 64 & 256 \\
Discount $\gamma$ & 0.99 & 0.99 & 0.99 \\
Buffer size & 50k & -- & 50k \\
Learning starts & 1000 & -- & 1000 \\
Train steps & 200k & 200k & 30k \\
\bottomrule
\end{tabular}
\end{table}

\subsection{Reproducibility}

The experiment pipeline stores trained models under \texttt{models/seed\_*}
and evaluation outputs under \texttt{results/seed\_*}. Each evaluation summary
is saved as JSON, while per-episode data are saved as CSV. This makes it
possible to recompute the aggregate tables without retraining the agents. The
final comparison tables are generated by averaging the per-seed metrics and
reporting both the mean and the sample standard deviation.

\section{Results}

\subsection{Training Metrics}

Table~\ref{tab:train} reports the number of completed training episodes,
collisions, successes, timeouts, and the collision rate for seed 0.

\begin{table}[htbp]
\caption{Training metrics for seed 0.}
\label{tab:train}
\centering
\begin{tabular}{lrrrrr}
\toprule
Algo & Episodes & Coll. & Succ. & Timeout & Coll. rate \\
\midrule
DQN & 5614 & 3811 & 1265 & 538 & 0.679 \\
PPO & 4784 & 1849 & 2345 & 590 & 0.386 \\
SAC & 425 & 252 & 0 & 173 & 0.593 \\
\bottomrule
\end{tabular}
\end{table}

PPO reached the highest number of successful episodes during training. DQN
also learned useful behavior, but with a higher collision rate. SAC did not
reach the goal during training under the selected budget.

\subsection{Deterministic Evaluation}

Table~\ref{tab:det} shows deterministic evaluation results averaged across
three seeds. PPO and DQN both achieved a mean success rate of 0.667. SAC timed
out in all evaluation episodes.

The aggregate mean hides the fact that the result is strongly seed-dependent.
Table~\ref{tab:seed_success} shows the deterministic success rate for each
individual seed. DQN and PPO both solve the fixed scene on seeds 0 and 1, but
both fail on seed 2. This indicates that the initialization and exploration
history have a large effect on the final policy in this small experiment.

\begin{table}[htbp]
\caption{Deterministic success rate per seed.}
\label{tab:seed_success}
\centering
\begin{tabular}{lrrr}
\toprule
Algorithm & Seed 0 & Seed 1 & Seed 2 \\
\midrule
DQN & 1.000 & 1.000 & 0.000 \\
PPO & 1.000 & 1.000 & 0.000 \\
SAC & 0.000 & 0.000 & 0.000 \\
\bottomrule
\end{tabular}
\end{table}

\begin{table}[htbp]
\caption{Deterministic evaluation, mean $\pm$ standard deviation.}
\label{tab:det}
\centering
\begin{tabular}{lrrrr}
\toprule
Algo & Success & Collision & Timeout & Steps succ. \\
\midrule
DQN & $0.667 \pm 0.577$ & $0.333 \pm 0.577$ & $0.000 \pm 0.000$ & $25.0 \pm 1.4$ \\
PPO & $0.667 \pm 0.577$ & $0.000 \pm 0.000$ & $0.333 \pm 0.577$ & $15.0 \pm 0.0$ \\
SAC & $0.000 \pm 0.000$ & $0.000 \pm 0.000$ & $1.000 \pm 0.000$ & -- \\
\bottomrule
\end{tabular}
\end{table}

The path quality comparison is shown in Table~\ref{tab:path}. DQN is close to
the A* grid optimum. PPO has a path ratio below 1 because it is not restricted
to axis-aligned grid moves and can move diagonally in continuous space.

\begin{table}[htbp]
\caption{Deterministic trajectory quality.}
\label{tab:path}
\centering
\begin{tabular}{lrrr}
\toprule
Algo & Path/A* & Avg reward & A* length \\
\midrule
DQN & $1.042 \pm 0.059$ & $86.7 \pm 138.8$ & 24.0 \\
PPO & $0.621 \pm 0.005$ & $118.3 \pm 82.2$ & 24.0 \\
SAC & -- & $-7.9 \pm 5.0$ & 24.0 \\
\bottomrule
\end{tabular}
\end{table}

\subsection{Stochastic Evaluation}

Table~\ref{tab:stoch} reports stochastic evaluation. PPO remains the strongest
continuous-control agent, while DQN keeps a comparable success rate but has
more collisions due to action sampling.

\begin{table}[htbp]
\caption{Stochastic evaluation, mean $\pm$ standard deviation.}
\label{tab:stoch}
\centering
\begin{tabular}{lrrrr}
\toprule
Algo & Success & Collision & Timeout & Steps succ. \\
\midrule
DQN & $0.589 \pm 0.517$ & $0.400 \pm 0.498$ & $0.011 \pm 0.019$ & $26.7 \pm 2.5$ \\
PPO & $0.633 \pm 0.549$ & $0.056 \pm 0.019$ & $0.311 \pm 0.539$ & $18.2 \pm 1.2$ \\
SAC & $0.000 \pm 0.000$ & $0.000 \pm 0.000$ & $1.000 \pm 0.000$ & -- \\
\bottomrule
\end{tabular}
\end{table}

\subsection{Trajectory Visualization}

Figure~\ref{fig:trajectories} shows representative deterministic trajectories
for seed 0. DQN follows the grid-like A* structure, PPO moves through smoother
continuous trajectories, and SAC fails to reach the target.

\begin{figure}[htbp]
\centering
\includegraphics[width=0.74\columnwidth]{../results/plots/trajectories_dqn.png}\\[-0.5ex]
\includegraphics[width=0.74\columnwidth]{../results/plots/trajectories_ppo.png}\\[-0.5ex]
\includegraphics[width=0.74\columnwidth]{../results/plots/trajectories_sac.png}
\caption{Deterministic trajectories for DQN, PPO, and SAC on seed 0. The green
line denotes the A* path, while blue lines denote agent trajectories.}
\label{fig:trajectories}
\end{figure}

\section{Discussion}

The results indicate that PPO is the best algorithm for this simplified
continuous 3D navigation task under the selected budget. It achieved no
deterministic evaluation collisions and produced the shortest successful
trajectories. DQN also solved the task, but its discrete action space causes
longer paths. This is expected because DQN can only move along one axis per
step.

SAC did not converge within 30,000 training steps. This does not necessarily
mean SAC is unsuitable for the task; rather, the selected budget was too small
for stable SAC learning in this setup. A more complete comparison would train
SAC for a substantially larger number of environment steps.

The high standard deviation in the success rate is caused by the small number
of seeds. Since the environment is fixed and the outcome is close to binary at
the seed level, three seeds are enough for a preliminary comparison but not for
a statistically strong conclusion.

The stochastic evaluation provides an additional view of policy robustness. A
deterministic policy can look strong if it has found one narrow trajectory, but
stochastic sampling reveals whether the learned action distribution also
assigns probability mass to safe actions around that trajectory. PPO remains
reasonably successful under stochastic evaluation, although its timeout rate
increases. DQN keeps a similar success rate but also produces more collisions,
which is expected because random deviations in a discrete grid can move the
agent directly into an obstacle or outside the safe corridor.

\section{Limitations}

The main limitation of the environment is its simplified physics. The drone is
modeled as a point mass with direct position control, so inertia, orientation,
motor dynamics, and realistic flight constraints are not represented. This
makes the task suitable for comparing RL algorithms, but it is not a full drone
simulator.

The observation is also intentionally compact. It contains the goal and only
the nearest obstacle, so the agent does not receive a full map of the world. In
the current scene with two obstacles this is sufficient, but in a more complex
environment the policy could fail because important obstacles may be hidden
from the observation vector.

Finally, the statistical evaluation uses only three seeds. The results are
useful for a course-scale experiment, but more seeds would be required for a
stronger empirical claim. Future work should use randomized start-goal pairs,
more obstacle layouts, and a longer SAC training budget.

\section{Conclusion}

This work implemented a reproducible 3D drone navigation benchmark and compared
DQN, PPO, and SAC on the same environment. PPO achieved the best overall
performance, combining high success rate, low collision rate, and short
continuous trajectories. DQN also solved the task but produced longer
grid-constrained paths. SAC did not learn a successful policy under the
selected training budget. Future work should include more seeds, randomized
start-goal pairs, richer obstacle observations, and longer SAC training.

\begin{thebibliography}{00}

\bibitem{madaan2020airsim}
R. Madaan, N. Gyde, S. Vemprala, M. Brown, K. Nagami, T. Taubner,
E. Cristofalo, D. Scaramuzza, M. Schwager, and A. Kapoor, ``AirSim Drone
Racing Lab,'' in \emph{Proceedings of the NeurIPS 2019 Competition and
Demonstration Track}, ser. Proceedings of Machine Learning Research, vol. 123,
PMLR, 2020, pp. 177--191.

\bibitem{mnih2015dqn}
V. Mnih et al., ``Human-level control through deep reinforcement learning,''
\emph{Nature}, vol. 518, no. 7540, pp. 529--533, 2015.

\bibitem{schulman2017ppo}
J. Schulman, F. Wolski, P. Dhariwal, A. Radford, and O. Klimov, ``Proximal
policy optimization algorithms,'' arXiv:1707.06347, 2017.

\bibitem{haarnoja2018sac}
T. Haarnoja, A. Zhou, P. Abbeel, and S. Levine, ``Soft actor-critic:
Off-policy maximum entropy deep reinforcement learning with a stochastic
actor,'' in \emph{Proc. ICML}, 2018.

\bibitem{stablebaselines3}
A. Raffin et al., ``Stable-Baselines3: Reliable reinforcement learning
implementations,'' \emph{Journal of Machine Learning Research}, vol. 22,
no. 268, pp. 1--8, 2021.

\bibitem{gymnasium}
M. Towers et al., ``Gymnasium,'' 2023. [Online]. Available:
\url{https://gymnasium.farama.org/}

\bibitem{hart1968astar}
P. E. Hart, N. J. Nilsson, and B. Raphael, ``A formal basis for the heuristic
determination of minimum cost paths,'' \emph{IEEE Transactions on Systems
Science and Cybernetics}, vol. 4, no. 2, pp. 100--107, 1968.

\end{thebibliography}

\end{document}
